\begin{abstract}
    Recent work by Garg~\etal~\cite{garg2024cifar10} demonstrated that a carefully engineered CNN can reach $94\%$ accuracy on \cifar{} in 3.29 seconds on a single GPU, using four training tricks: frozen patch-whitening initialisation, Dirac (identity) initialisation, a ScaleBias output layer with the Lookahead optimiser, and alternating horizontal flip.
    We ask whether these tricks transfer to transformer architectures.
    We train a small Vision Transformer (ViT, $\sim$4M parameters) on \cifar{} under two regimes: a weak recipe (100 epochs, basic augmentation, $61.7\%$ baseline) and a strong recipe (300 epochs, RandAugment, Mixup/CutMix, DropPath, $81.4\%$ baseline).
    Under the weak recipe, representation-quality tricks transfer positively (whitening $+1.3$\pp{}, alternating flip $+0.8$\pp{}) while architecture-specific tricks are harmful (Dirac $-7.9$\pp{}, ScaleBias+Lookahead $-16.0$\pp{}).
    Under the strong recipe the picture partially reverses: whitening flips to $-4.4$\pp{} and soft Dirac init ($\alpha{=}0.3$) becomes the only net-positive trick ($+0.4$\pp{}).
    We attribute the whitening reversal to a frozen-embedding/augmentation conflict: filters computed from clean images are mismatched to the heavily transformed inputs of Mixup and CutMix.
    A controlled CNN masked-reconstruction experiment confirms that freezing itself is harmless---only the statistics mismatch matters.
    These findings reveal that the transferability of CNN training tricks depends on the \emph{trick~$\times$~recipe} interaction, not just on architectural compatibility.
\end{abstract}
